% Quelle: klausur PZR1.png
% Art: Klausur/Klausurvorbereitung
% Themen: RSA, Binärdarstellung, Eulersche Phi-Funktion, Inklusion-Exklusion, Euklidischer Algorithmus, Chinesischer Restsatz, Modulo-11-Code, Kontrollmatrix, Syndrom, Decodierung
\section{Klausur vom 29.07.2026}
% Quelle ist ein Foto des Aufgabenblatts (Screenshot eines Messenger-Bildes); nur Aufgabenstellungen, keine Lösungen vorhanden.
% Kopfzeile des Blatts: Name: \quad Matrikelnr.:
% Punktzahlen je Aufgabe sind auf dem Blatt nicht angegeben.

\subsection{Aufgaben}

\begin{aufgabe}[Aufgabe 1]
% Punkte: nicht angegeben
Alice will die Nachricht $w = 1511$ (GO) verschlüsseln und als verschlüsselte Nachricht $c$ (Codewort) an Bob senden. Sie verabreden als Modul die Zahl $m = 43 \cdot 57 = 2451$ zu wählen und den öffentlichen Schlüssel $e = 221$ zu verwenden.
\begin{enumerate}[label=\alph*)]
  \item Schreiben Sie $e$ im Binärsystem.
  \item Mit welcher Formel verschlüsselt Alice das Wort $w$? Berechnen Sie das Codewort $c$.
  \item Bob besitzt den privaten Schlüssel $d$. Wie (Formel) entschlüsselt Bob das Codewort $c$?
  \item Berechnen Sie $\varphi(m)$.
  \item Oscar erhält Kenntnis von $m$ und $e$. Welche Formeln muss er nutzen um den privaten Schlüssel $d$ zu berechnen? Berechnen Sie $d$.
\end{enumerate}
\end{aufgabe}

\begin{aufgabe}[Aufgabe 2]
% Punkte: nicht angegeben
Wie viele positive ganze Zahlen $\leq 720$ sind weder durch $3$, $4$ oder $5$ teilbar. Verwenden Sie die Symbole $A_i$, $\alpha_i$ aus der Vorlesung. (Inklusion -- Exklusion)
\end{aufgabe}

\begin{aufgabe}[Aufgabe 3]
% Punkte: nicht angegeben
Bestimmen Sie mit Hilfe des Euklidischen Algorithmus
\begin{enumerate}[label=\alph*)]
  \item den größten gemeinsamen Teiler von $m = 17$ und $n = 15$,
  \item eine ganzzahlige Lösung der Gleichung $17x - 15y = 1$.
\end{enumerate}
\end{aufgabe}

\begin{aufgabe}[Aufgabe 4]
% Punkte: nicht angegeben
Gegeben ist das System simultaner Kongruenzen (Chinesischer Restsatz)
\begin{align*}
  z &\equiv 8 \pmod{17}\\
  z &\equiv 5 \pmod{15}
\end{align*}
\begin{enumerate}[label=\alph*)]
  \item Bestimmen Sie eine ganze Zahl $z$ mit Hilfe des Ergebnisses aus Aufgabe 3.
  \item Bestimmen Sie sämtliche Lösungen dieses Systems, insbesondere die kleinste positive Lösung.
\end{enumerate}
\end{aufgabe}

\begin{aufgabe}[Aufgabe 5]
% Punkte: nicht angegeben
Die Parameter eines Modulo-11 Codes sind $[10, 8, 3]_{11}$.\\
Vorgelegt ist das Klartextwort $a = 1\,1\,5\,0\,0\,0\,0\,5$.
\begin{enumerate}[label=\alph*)]
  \item Wie lautet die Kontrollmatrix $H$?
  \item Berechnen Sie die Kontrollziffern $c_9$, $c_{10}$ und damit das zugehörige Codewort $c$.
\end{enumerate}
Empfangen wurden die Wörter
\begin{align*}
  r_1 &= 1\,1\,5\,0\,0\,0\,0\,5\,1\,1\\
  r_2 &= 1\,1\,5\,0\,0\,0\,0\,7\,3\,3
\end{align*}
\begin{enumerate}[label=\alph*), start=3]
  \item Berechnen Sie die Syndrome $S(r_1)$, $S(r_2)$?
  \item Welches der beiden Wörter $r_1$, $r_2$ ist nicht dekodierbar und warum nicht?
  \item Dekodieren Sie das andere Wort und geben Sie das Codewort $c$ an.
  \item Zeigen Sie, dass $c$ tatsächlich ein Codewort ist.
\end{enumerate}
\end{aufgabe}

\subsection{Bewertung}
Punktzahl: \qquad Note:

\noindent Bewertungsschema:

\begin{tabular}{*{10}{c}}
\toprule
29/1.0 & 28/1.3 & 26/1.7 & 25/2.0 & 24/2.3 & 22/2.7 & 21/3.0 & 19/3.3 & 17/3.7 & 15/4.0\\
\bottomrule
\end{tabular}
